\begin{examplebox}
	\textbf{\textcolor{CExample}{例 1（基础）}}

	\textbf{\textcolor{CExample}{题目：}}
	已知椭圆 $C: \frac{x^2}{25}+\frac{y^2}{16}=1$ 的左、右焦点分别为 $F_1,F_2$，点 $P$ 在 $C$ 上且 $P$ 的横坐标为 $5$。求 $|PF_1|$ 和 $|PF_2|$。

	\textbf{\textcolor{CExample}{解题步骤：}}
	\begin{description}[style=nextline, leftmargin=2.8em, labelsep=0.5em, itemsep=0.45em]
		\item[\textbf{第一步（识别参数）：}] 由椭圆标准方程得 $a=5$，$b=4$，$c=\sqrt{a^2-b^2}=3$，离心率 $e=\frac{c}{a}=\frac{3}{5}$。
			\par\textbf{理由：} 椭圆方程化为标准式，对比系数。
		\item[\textbf{第二步（代入公式）：}] 点 $P$ 的横坐标 $x_0=5$，代入焦半径公式：
			\[
				|PF_1| = a + e x_0 = 5 + \frac{3}{5}\times 5 = 8,
				\qquad
				|PF_2| = a - e x_0 = 5 - \frac{3}{5}\times 5 = 2.
			\]
			\par\textbf{理由：} 椭圆焦半径公式 $|PF_1|=a+ex_0,\ |PF_2|=a-ex_0$。
		\item[\textbf{第三步（结论）：}] 所以 $|PF_1|=8$，$|PF_2|=2$。
			\par\textbf{理由：} 计算结果即为所求。
	\end{description}

	\textbf{\textcolor{CExample}{关键结论：}}
	\textbf{$|PF_1| = 8$，$|PF_2| = 2$。}

	\medskip
	\textbf{\textcolor{CExample}{例 2（稍变形）}}

	\textbf{\textcolor{CExample}{题目：}}
	已知椭圆 $C: \frac{x^2}{25}+\frac{y^2}{16}=1$，点 $P$ 在 $C$ 上，且 $|PF_2| = 4$。求点 $P$ 的横坐标 $x_0$。

	\textbf{\textcolor{CExample}{解题步骤：}}
	\begin{description}[style=nextline, leftmargin=2.8em, labelsep=0.5em, itemsep=0.45em]
		\item[\textbf{第一步（建立方程）：}] 由焦半径公式 $|PF_2| = a - e x_0$，代入 $a=5$，$e=\frac{3}{5}$，$|PF_2|=4$，得 $5 - \frac{3}{5}x_0 = 4$。
			\par\textbf{理由：} 焦半径公式的直接使用。
		\item[\textbf{第二步（解方程）：}] 整理得 $\frac{3}{5}x_0 = 1$，解得 $x_0 = \frac{5}{3}$。
			\par\textbf{理由：} 一元一次方程求解。
		\item[\textbf{第三步（检验范围）：}] 因为 $x_0 = \frac{5}{3} \in [-5,5]$，所以存在点 $P$ 满足条件。
			\par\textbf{理由：} 椭圆上点的横坐标范围是 $[-a,a]$。
	\end{description}

	\textbf{\textcolor{CExample}{关键结论：}}
	\textbf{点 $P$ 的横坐标为 $\displaystyle \frac{5}{3}$。}

	\medskip
	\textbf{\textcolor{CExample}{例 3（综合应用）}}

	\textbf{\textcolor{CExample}{题目：}}
	已知椭圆 $C: \frac{x^2}{25}+\frac{y^2}{16}=1$，点 $P$ 在椭圆上。求 $|PF_1|\cdot|PF_2|$ 的最大值。

	\textbf{\textcolor{CExample}{解题步骤：}}
	\begin{description}[style=nextline, leftmargin=2.8em, labelsep=0.5em, itemsep=0.45em]
		\item[\textbf{第一步（乘积表达）：}] 由焦半径公式得
			\[
				|PF_1|\cdot|PF_2| = (a+ex_0)(a-ex_0) = a^2 - e^2 x_0^2 = 25 - \frac{9}{25}x_0^2.
			\]
			\par\textbf{理由：} 平方差公式，代入 $a=5,\ e=\frac{3}{5}$。
		\item[\textbf{第二步（范围条件）：}] 点 $P$ 在椭圆上，故 $x_0\in[-5,5]$。
			\par\textbf{理由：} 椭圆上点的横坐标范围 $[-a,a]$。
		\item[\textbf{第三步（求最值）：}] 函数 $f(x_0)=25-\frac{9}{25}x_0^2$ 在 $[-5,5]$ 上是偶函数，最大值出现在 $x_0=0$ 处，最大值为 $25$。
			\par\textbf{理由：} 二次函数开口向下，顶点处取最大值。
	\end{description}

	\textbf{\textcolor{CExample}{关键结论：}}
	\textbf{$|PF_1|\cdot|PF_2|$ 的最大值为 $25$。}


\medskip
\textbf{\textcolor{CExample}{例 4（从焦半径之差反求坐标）}}

\textbf{\textcolor{CExample}{题目：}}
椭圆 $\frac{x^2}{25}+\frac{y^2}{16}=1$ 上一点 $P$ 满足
$|PF_1|-|PF_2|=\frac{12}{5}$。求 $P$ 的横坐标。

\textbf{\textcolor{CExample}{解法：}}
从平方差法得到的关键关系 $|PF_1|-|PF_2|=2ex_0$ 出发。
这里 $a=5$，$c=3$，$e=\frac35$，故
\[
 \frac{12}{5}=2\cdot\frac35\,x_0
 \quad\Longrightarrow\quad \boxed{x_0=2}.
\]
检验 $2\in[-5,5]$，椭圆上确有这样的点。
\textbf{识别提示：} 已知“两条焦半径的差”，立刻想到 $2ex_0$，不必设出 $y_0$。

\medskip
\textbf{\textcolor{CExample}{例 5（焦点改在 $y$ 轴）}}

\textbf{\textcolor{CExample}{题目：}}
椭圆 $\frac{x^2}{9}+\frac{y^2}{25}=1$ 的下、上焦点分别为 $F_1,F_2$。
设点 $P$ 在椭圆上，纵坐标 $y_0=\frac52$，求两条焦半径。

\textbf{\textcolor{CExample}{解法：}}
$a=5$，$b=3$，$c=4$，$e=\frac45$。
焦点在 $y$ 轴，所以用\emph{纵坐标}：
\[
 |PF_1|=5+\frac45\cdot\frac52=\boxed{7},\qquad
 |PF_2|=5-\frac45\cdot\frac52=\boxed{3}.
\]
两条距离之和 $7+3=10=2a$，与第一定义相符。
\textbf{识别提示：} 焦点轴是哪条，就用哪条轴上的坐标；下加上减。
\end{examplebox}
